\chapter{Metodología de investigación}
\label{chap:methodology}

\section{Paradigma de investigación y tipo de estudio}

This thesis adopts a design-science paradigm because its central research output is an artifact: a coordination framework with explicit constructs, decision rules, an architectural instantiation, and an evaluation method. The study combines exploratory and analytical work in the literature review, formal and architectural design in the construction phase, and quantitative and qualitative analysis in the planned evaluation. Quantitative measures assess classification, task execution, latency, and overhead; qualitative trace analysis examines whether authorization and refusal decisions are intelligible and supported by recorded evidence.

\section{Proceso de investigación}

The research process contains five linked stages. First, a critical literature review identifies the strengths and unresolved limitations of classical multi-agent coordination, symbolic reasoning, LLM agents, and interoperability protocols. Second, the research gap is translated into requirements for openness, interoperability, bounded authority, traceability, and explicit refusal. Third, these requirements are formalized as a protocol-independent framework. Fourth, the framework is mapped to a Python \texttt{CoordinationAgent} design over SDK middleware that uses A2A as its interoperability backbone. Fifth, controlled scenarios compare the complete design with baselines and ablations. Each stage produces evidence for one or more specific objectives rather than serving as an independent implementation activity.

\begin{table}[htbp]
\centering
\caption{Alineación de objetivos, métodos y evidencia.}
\label{tab:method-alignment}
\begin{tabular}{p{0.12\textwidth}p{0.38\textwidth}p{0.38\textwidth}}
\toprule
\textbf{Objetivo} & \textbf{Método} & \textbf{Evidencia} \\
\midrule
1 & Síntesis temática crítica de la literatura & Criterios de comparación y brecha explícita de investigación \\
2 & Modelado formal y conceptual & Definiciones, reglas, predicados de viabilidad y propiedades declaradas \\
3 & Diseño arquitectónico y mapeo de protocolo & Componentes, registros, flujos de trabajo y trazabilidad teoría-implementación \\
4 & Evaluación controlada por escenarios & Comparaciones con líneas base, ablaciones, métricas y trazas de ejecución \\
5 & Discusión analítica & Análisis de alcance, casos de fallo, limitaciones y amenazas a la validez \\
\bottomrule
\end{tabular}
\end{table}

\section{Unidad de análisis y diseño de escenarios}

The primary unit of analysis is a coordination attempt: one problem request, one admitted registry snapshot, the candidate plan or refusal produced for that request, and the resulting execution trace when a plan is authorized. Scenarios are constructed by varying four factors: capability coverage, artifact compatibility, authority constraints, and runtime availability. This factorial view prevents evaluation from relying only on successful demonstrations.

The scenario suite must contain feasible direct-delegation cases, feasible decomposed cases, cases requiring an approved bounded linguistic transformation, and infeasible cases caused by missing capability, incompatible modality, cyclic dependency, insufficient authority, or unavailable agents. Each scenario includes a reference feasibility label and the minimum justification required for that label. Reference labels should be reviewed independently before system outputs are examined.

\section{Líneas base y ablaciones}

Three configurations provide the minimum comparison. The \emph{hybrid CoordinationAgent} uses \texttt{CoordinationSdk} as the only communication boundary, may invoke Lingo-backed SDK capabilities for structured interpretation and candidate decomposition, and always applies symbolic feasibility checks before dispatch. The \emph{CoordinationAgent without symbolic feasibility} variant receives the same request and registry descriptions but accepts or rejects its own proposed plan without the deterministic feasibility layer. The \emph{rule-only CoordinationAgent} variant receives a structured request and applies the same explicit predicates without LLM interpretation. The first baseline tests the contribution of symbolic authorization; the second tests the contribution of linguistic mediation. Mixed-agent scenarios must include remote A2A agents, local Python agents, and Lingo-backed linguistic agents so that all specializations are evaluated through the same SDK boundary.

Additional ablations remove one mechanism at a time: contract validation, explicit authority checks, trace recording, or bounded-auxiliary rules. These ablations are evaluated only where the removed mechanism is relevant to the scenario. They are not treated as independent general-purpose systems.

\section{Dimensiones y métricas de evaluación}

\begin{table}[htbp]
\centering
\caption{Dimensiones de evaluación, indicadores y evidencia.}
\label{tab:method-metrics}
\begin{tabular}{p{0.24\textwidth}p{0.30\textwidth}p{0.34\textwidth}}
\toprule
\textbf{Dimensión} & \textbf{Indicador} & \textbf{Métrica o evidencia} \\
\midrule
Juicio de viabilidad & Autorización o rechazo correcto & Exactitud, precisión, exhaustividad y F1 para la clasificación de planes viables \\
Validez del plan & Detección de candidatos inválidos & Tasa de rechazo de planes inválidos por categoría de fallo \\
Resultado de ejecución & Artefactos requeridos completados & Tasa de éxito de tareas y tasa de artefactos válidos por contrato \\
Trazabilidad & Evidencia requerida de decisión registrada & Completitud de campos de traza y tasa de decisiones reconstruibles \\
Eficiencia & Esfuerzo de coordinación & Latencia extremo a extremo, llamadas al planificador, mensajes y uso de tokens \\
Robustez & Estabilidad bajo variación controlada & Desempeño ante paráfrasis, pérdida de agentes y perturbación de metadatos \\
\bottomrule
\end{tabular}
\end{table}

Outcome metrics and process metrics are reported separately. A correct final artifact does not erase an unauthorized action or an invalid intermediate plan. Conversely, a justified refusal is a correct outcome for an infeasible request and must not be counted as task failure.

\section{Procedimiento experimental}

For each scenario, every applicable configuration receives semantically equivalent inputs and an identical admitted-agent snapshot. Stochastic configurations are executed repeatedly with recorded model, decoding parameters, prompt version, and random seed where supported. Each run stores the interpreted request, candidate plan, predicate outcomes, authorization decision, coordination-agent decisions, SDK method calls, agent kind, transport or runtime path, messages, status transitions, artifacts, validation results, latency, and resource use. Aggregate results are accompanied by per-category error analysis. The implementation version, coordination-agent version, SDK version, protocol version, hardware, and software environment must be fixed and reported before measurements are collected.

\section{Estrategia de análisis}

Feasibility decisions are compared with reference labels through a confusion matrix and category-level error rates. Paired scenario results are used for comparisons because configurations face the same requests. Descriptive statistics include central tendency and dispersion rather than averages alone. Where the sample size and distribution justify inference, paired statistical tests and effect sizes are reported. Qualitative analysis classifies false acceptance, false refusal, invalid decomposition, unsupported auxiliary synthesis, authority violation, artifact mismatch, and incomplete evidence.

\section{Validez y reproducibilidad}

Construct validity depends on whether the metrics capture justified coordination rather than surface task completion. This threat is mitigated by separating plan validity, execution success, communication-boundary compliance, coordination-ownership compliance, and trace completeness. Internal validity may be affected by unequal prompts, model nondeterminism, or baseline implementation quality; shared inputs, repeated runs, versioned configurations, and paired comparisons reduce these effects. External validity is limited by the selected scenarios, admitted agents, SDK adapters, and A2A-backed implementation choices. The theoretical framework is protocol-independent, but empirical findings may not generalize to decentralized authorization, physical robots, adversarial registries, or other protocols. Conclusion validity depends on scenario count and variance and must therefore be assessed after data collection.

Reproducibility requires publication of scenario specifications, reference labels, registry snapshots, configuration files, prompts, raw traces, analysis scripts, and exact dependency versions. Sensitive credentials and private agent endpoints are excluded, but their interfaces and response fixtures should be represented by reproducible substitutes.
